\section{Discussion and future work} \label{sec:discussion}
In this article, we have introduced Soft Mapper, a distributional and smoother version of the standard Mapper graph, with provable convergence guarantees using persistence-based losses and risks. Our case study in this article was finding an optimal filter function, among a parameterized family of functions, in order to construct regular Mapper graphs incorporating an optimized and maximal amount of topological information. 
We then produced examples of such optimization processes on real 3D shape and single-cell RNA sequencing data, for which we were able to obtain structured Mapper representations in an unsupervised way.
%To achieve this, we used persistent homology, to produce a topological loss, and a specific Soft Mapper that can be seen a smooth relaxation of the regular Mapper. We then used this method to produce data representations based on linear filters. 
These representations, especially for the single cell RNA-sequencing data, are not meant to represent novel or state of the art data representations in their respective research domains, but as a proof of concept of the practical benefit of our method.
Moreover, our construction is not limited to the choice of a linear family of filter functions or to the filter optimization setting as a whole. Possible future work includes optimizing non-linear filter functions, based on neural networks or kernel methods, and studying Soft Mappers based on different cover assignment schemes, like the Gaussian cover assignment scheme defined in Appendix \ref{apdx:gaussian}. 

\section{Acknowledgments}

The research was supported by three grants from Agence Nationale de la Recherche: GeoDSIC ANR-22-CE40-0007, ANR JCJC TopModel ANR-23-CE23-0014 and AI4scMed, France 2030 ANR-22-PESN-000. 

